\chapter{Diskussion}
\label{cha: diskussion}

Dieses Kapitel interpretiert die in Kapitel~\ref{cha:ergebnisse} berichteten Ergebnisse: Zunächst
werden die Forschungsfragen beantwortet (Abschnitt~\ref{sec:beantwortung-forschungsfragen}) und die
Befunde in den Forschungsstand eingeordnet (Abschnitt~\ref{sec:einordnung-forschungsstand}),
anschließend werden die Grenzen der Untersuchung dargelegt (Abschnitt~\ref{sec:limitationen}) und ein
Ausblick auf weiterführende Arbeiten gegeben (Abschnitt~\ref{sec:ausblick}).

\section{Beantwortung der Forschungsfragen}
\label{sec:beantwortung-forschungsfragen}

Die in Abschnitt~\ref{sec:forschungsfrage} formulierten Teilfragen werden im Folgenden auf Grundlage
der qualitativen Kodierung (Abschnitt~\ref{sec:datenanalyse}), der Survey-Auswertung
(Abschnitt~\ref{sec:fragebogen-experiment}) sowie der in Kapitel~\ref{cha:ergebnisse} berichteten
Befunde beantwortet, bevor abschließend die übergeordnete Forschungsfrage zusammengeführt wird.

\textbf{TF1 -- Auslagerung und ungeprüfte Übernahme.} Der Reflexionsbot verzichtet bewusst auf
direkte Antworten und passt die Tiefe seiner Rückfragen an die Qualität der studentischen
Argumentation an (adaptive Tiefe, Abschnitt~\ref{sec:konzept-reflexionsbot}). Der deutlichste
Unterschied zeigt sich im Anteil vollständiger oder teilweiser Auslagerung der Bearbeitung an die
\acrshort{ki}: Bei Team AI sind 80\,\% der kodierten Fälle (12 von 15) nicht eindeutig studentisch,
beim Reflexionsbot sind es 50\,\% (8 von 16, vgl. Abschnitt~\ref{sec:ergebnisse}). Das
Rückfrage-Prinzip geht mit einer geringeren Auslagerungsquote einher, ohne die Auslagerung vollständig
zu verhindern, wie die Umgehungsversuche einzelner Personen zeigen (vgl. Abschnitt~\ref{sec:illustration-kategoriensystem}).
Zugleich ist der Abstand teilweise durch die Aufgabenvorgabe für Team AI mitbedingt und daher als
Indiz, nicht als exakte Effektgröße zu lesen (vgl. Abschnitt~\ref{sec:limitationen}). Ein Teil des
Unterschieds könnte zudem darauf beruhen, dass der Reflexionsbot keine direkt kopierbare Lösung
ausgibt, sodass die geringere Auslagerung nicht allein auf die didaktische Gestaltung, sondern auch
auf das Fehlen einer Kopiervorlage zurückgehen kann. In der
Generic-AI-Gruppe ist der Anteil ungeprüft übernommener Vorschläge damit höher als beim Reflexionsbot.
Gleichzeitig zeigt Abschnitt~\ref{sec:selbstauskunft-verhalten}, dass die Selbstauskunft im Survey
diesen Unterschied allein nicht zuverlässig abbildet: Personen, die nachweislich nur kopiert hatten,
bewerteten das jeweilige System teils trotzdem als hilfreich. Vertiefte Auseinandersetzung entsteht
also nicht automatisch durch das bloße Vorhandensein von Rückfragen, sondern hängt von der
Bereitschaft der Studierenden ab, sich auf den Dialog einzulassen.

\textbf{TF2 -- Reflexionstiefe und Prozessverständnis.} Beim selbsteingeschätzten Prozessverständnis
zeigt sich ein differenziertes Bild: Bei Items zum reinen Beschreiben sinkt die Selbsteinschätzung in
beiden Gruppen ähnlich stark. Bei Items zum Begründen und Abwägen bleibt die Reflexionsbot-Gruppe
stabil oder verbessert sich leicht, während sie in der Generic-AI-Gruppe durchgehend weiter sinkt
(vgl. Abschnitt~\ref{sec:ergebnisse}). Umgekehrt schätzt sich die Generic-AI-Gruppe bei mehreren
direkten Reflexions-Items des Post-Surveys (etwa \glqq überlegt, warum die eigenen Entscheidungen
angemessen waren\grqq{} oder \glqq mehrere Optionen verglichen\grqq{}) sogar etwas höher ein
(\hyperref[app:statistik]{Anhang~A.6}). Das selbstberichtete Bild ist damit insgesamt uneinheitlich und stützt keinen klaren
Systemvorteil.

In der tatsächlich gezeigten Reflexionstiefe der studentisch verfassten Bearbeitungen fällt das
Ergebnis deutlicher aus, und zwar gegen die Ausgangsannahme dieser Arbeit: Die wenigen Team-AI-Fälle
liegen auf Stufe~2, die Reflexionsbot-Fälle überwiegend auf den Stufen~0 und~1 mit einem Einzelfall
auf Stufe~3 (Abb.~\ref{fig:ergebnisse-stufen}). Deskriptiv liegt die Vergleichsgruppe damit sowohl im
Mittelwert (1.33 gegenüber 0.75) als auch im Median (2.0 gegenüber 0.5) höher. Auf dem Konstrukt, das
der Reflexionsbot gezielt fördern soll, schneidet er in dieser Stichprobe also nicht besser ab. Bei
drei gegenüber acht auswertbaren Fällen ist dieser Unterschied allerdings nicht belastbar. Bereits ein
einzelner Fall verschiebt den Mittelwert der kleineren Gruppe um mehr als 0,3 Stufen, und die drei
Team-AI-Fälle sind der Rest einer Gruppe, in der 80\,\% der Bearbeitungen als nicht eindeutig
studentisch ausgeschieden sind. Es ist daher ebenso möglich, dass hier eine positive Selektion der
wenigen verbliebenen Fälle sichtbar wird wie ein Nachteil des Reflexionsbots. Als belastbarer Befund
bleibt festzuhalten, dass sich ein Vorteil des Reflexionsbots auf dem Zielkonstrukt in diesen Daten
nicht zeigt.

Am objektiven Wissenstest zeigt sich zwischen den Systemgruppen ebenfalls kein Vorteil des
Reflexionsbots, wobei der Deckeneffekt hier jede Richtungsaussage verbietet
(vgl. Abschnitt~\ref{sec:ergebnisse}). Auch die Aufteilung nach dem Bearbeitungsverhalten trägt
weniger, als sie auf den ersten Blick verspricht: Die eigenständig arbeitenden Personen lagen bereits
im Prä-Test durchgehend bei der vollen Punktzahl, sodass ihr besseres Post-Ergebnis überwiegend eine
vorbestehende Leistungsdifferenz abbildet und keinen Lernzuwachs. Der Zusammenhang zwischen
eigenständiger Bearbeitung und Testergebnis ist in diesen Daten daher als Selektionseffekt zu lesen.
Die inhaltlich tragfähigere Beobachtung ist eine andere: Wer die Aufgabe an die \acrshort{ki}
auslagerte, erzeugte Texte, die sich einer Reflexionskodierung von vornherein entziehen. Nicht das
Testergebnis, sondern die Bearbeitungshaltung entscheidet darüber, ob überhaupt etwas entsteht, an dem
sich Lernen zeigen könnte.

\textbf{TF3 -- Eingeschätzte Kompetenzentwicklung.} Die vier hierfür einschlägigen Survey-Items sind
in Abschnitt~\ref{sec:ergebnisse} berichtet und in Abbildung~\ref{fig:survey-vergleich}
gegenübergestellt. Sie ergeben ein gegenläufiges Bild. Die wahrgenommene Kompetenzentwicklung fällt
beim Reflexionsbot leicht positiver aus (M=3.10 gegenüber M=2.80), und das System übertraf die
Erwartungen der Studierenden deutlich stärker (M=3.40 gegenüber M=2.20). Zugleich ist die Bereitschaft
zur erneuten Nutzung geringer (M=2.50 gegenüber M=3.20) und der Wunsch nach dem jeweils anderen System
größer (M=3.30 gegenüber M=2.20).

Für das Auseinanderfallen von positiverer Einschätzung und geringerer Nutzungspräferenz liegt die
Deutung nahe, der Reflexionsbot werde als anstrengender, aber lohnender wahrgenommen. Die erhobene
Skala zur kognitiven Last stützt diese Deutung jedoch nicht eindeutig: Bei den Teilaufgaben liegt die
Reflexionsbot-Gruppe geringfügig höher, bei der zusammenfassenden Einschätzung des gesamten kognitiven
Aufwands dagegen niedriger als die Vergleichsgruppe (M=3.00 gegenüber M=3.20,
vgl. Abschnitt~\ref{sec:ergebnisse}). Die Anstrengungsdeutung bleibt damit eine plausible, aber
unbelegte Erklärung.

Zu berücksichtigen ist ferner, dass das Item \glqq war hilfreicher als erwartet\grqq{} keinen
absoluten Nutzen misst, sondern einen Abgleich mit einer Erwartung. Die Vergleichsgruppe arbeitete mit
Systemen, die sie nach eigener Angabe im Prä-Survey bereits regelmäßig nutzt. Eine positive
Überraschung ist dort strukturell unwahrscheinlicher als bei einem unbekannten System. Der deutliche
Abstand bei diesem Item kann daher ebenso gut die Neuheit des Reflexionsbots abbilden wie seine
didaktische Gestaltung.

\textbf{TF3 -- Wahrgenommene Stärken und Grenzen.} Als Stärke des Reflexionsbots zeigen die Daten eine
stärkere Anregung zu unabhängigem Denken \emph{als erwartet}, ein deutlicheres Übertreffen der
Erwartungen und ein stabileres selbsteingeschätztes Verständnis für Begründung und Abwägung. Der
Zusatz \glqq als erwartet\grqq{} ist dabei Teil des Items und nicht wegzulassen: Gemessen wurde ein
Erwartungsabgleich, nicht das Ausmaß unabhängigen Denkens selbst
(\hyperref[app:statistik]{Anhang~A.6}). Als Grenze zeigt sich eine breitere Streuung der
Nützlichkeitsbewertung zwischen \glqq sehr nützlich\grqq{} und \glqq nicht nützlich\grqq{}, die zur
stark rechtsschiefen Verteilung der Reflexionsstufen mit einem Einzelfall auf Stufe~3 passt
(Abschnitt~\ref{sec:ergebnisse}), sowie eine geringere Bereitschaft zur erneuten Nutzung. Für das
generische System gilt umgekehrt eine höhere wahrgenommene Convenience und Nutzungspräferenz, zugleich
aber ein deutlich höherer Anteil vollständiger KI-Auslagerung in den tatsächlichen Bearbeitungen.

\textbf{Zusammenführung zur zentralen Forschungsfrage.} Aus den drei Teilfragen ergibt sich ein
differenziertes und in Teilen gegenläufiges Bild. In den vorliegenden Daten steht ein didaktisch
gestalteter Reflexionsbot vor allem dort mit einer vertieften Auseinandersetzung in Verbindung, wo
sich Studierende auf den sokratischen Dialog einlassen: erkennbar an einer geringeren
KI-Auslagerungsquote, einem stabileren selbsteingeschätzten Prozessverständnis bei höherstufigen
Kompetenzen und einer positiveren Einschätzung der eigenen Kompetenzentwicklung im Vergleich zu einem
generischen System. Auf dem eigentlichen Zielkonstrukt, der in den Texten gezeigten Reflexionstiefe,
zeigt sich dagegen kein Vorteil; deskriptiv liegt die Vergleichsgruppe dort sogar höher, bei einer
Fallzahl, die keine Richtungsaussage trägt. Am objektiven Wissenstest lässt der Deckeneffekt
zwischen den Systemgruppen keine Aussage zu (Abschnitt~\ref{sec:ergebnisse}), und der Zusammenhang
mit eigenständiger Bearbeitung erweist sich bei Betrachtung der Prä-Werte als Selektionseffekt. Der
Reflexionsbot wirkt zudem
nicht automatisch und nicht für alle Studierenden gleichermaßen. Ein erheblicher Teil der Reflexionsbot-Gruppe vermied
die inhaltliche Auseinandersetzung dennoch weitgehend, wich Rückfragen aus oder versuchte, das System
zu direkten Antworten zu bewegen. Die entscheidende Voraussetzung für das beobachtete Muster ist
damit nicht das Werkzeug allein, sondern die Interaktionsbereitschaft der Studierenden, die durch das
Werkzeugdesign begünstigt, aber nicht erzwungen werden kann. Diese Deutung ist allerdings explorativ
und in Folgestudien gezielt zu prüfen.

\section{Einordnung in den Forschungsstand}
\label{sec:einordnung-forschungsstand}

Die vorliegenden Befunde lassen sich an mehreren Stellen mit der in Kapitel~\ref{chap:theogrundlage}
aufgearbeiteten Literatur in Verbindung setzen und ergänzen diese um empirische Evidenz aus einem
Themenfeld, das dort bislang kaum untersucht war.

TreeInstruct \cite{karguptaInstructNotAssist2024} zeigte, dass sokratisches Nachfragen Studierende
erfolgreich beim eigenständigen Auffinden einzelner Programmierfehler unterstützen kann, also bei
Aufgaben mit einer eindeutig richtigen Lösung. Die vorliegende Arbeit überträgt dieses Prinzip auf
Prozessentscheidungen ohne eindeutig richtige Lösung und deutet darauf hin, dass das Prinzip auch dort
greifen könnte: Der Anteil vollständig oder teilweise an die KI ausgelagerter Bearbeitungen fällt
geringer aus (Abschnitt~\ref{sec:beantwortung-forschungsfragen}). Anders als bei TreeInstruct, dessen
Fragebaum den Wissensstand der Studierenden aktiv modelliert, verlässt sich der hier entwickelte
Reflexionsbot auf eine einfachere, an der Qualität der jeweils letzten Antwort orientierte adaptive
Tiefe (Abschnitt~\ref{sec:konzept-reflexionsbot}). Dass Studierende diese vereinfachte Steuerung
dennoch gezielt umgehen konnten (Abschnitt~\ref{sec:ausblick}), deutet darauf hin, dass eine
differenziertere Modellierung des Antwortverhaltens, wie sie TreeInstruct für Programmierfehler
einsetzt, auch für offene Prozessentscheidungen einen lohnenden Ansatzpunkt bieten könnte.

\glqq Reflect Me\grqq{} \cite{mollerReflectMeTextgenerative2025} bietet allgemeine, domänenoffene
Reflexionsimpulse in Form vorformulierter Prompts, ohne deren Wirksamkeit empirisch zu prüfen. Welche
Fragetypen tatsächlich zu tieferer Reflexion führen, bleibt dort offen. Die vorliegende Arbeit
liefert hierzu einen Anhaltspunkt:
Gerade die Szenario-Spezifität, also Rückfragen, die konkrete Randbedingungen und Stakeholder eines
bestimmten Dilemmas benennen, statt allgemein gehaltene Reflexionsfragen zu stellen, erwies sich als
Voraussetzung dafür, dass die adaptive Tiefe überhaupt präzise auf einzelne Argumentationslücken
reagieren konnte (Abschnitt~\ref{sec:konzept-reflexionsbot}). Für Themenfelder mit konkreten,
fachlich geprägten Entscheidungssituationen wie dem Software-Prozessmanagement spricht dies für enger
gefasste, kontextspezifische Reflexionswerkzeuge statt allgemeiner Prompt-Sammlungen.

Eine Übersichtsarbeit zum kognitiven Einfluss generativer KI-Werkzeuge wertet 17 Studien aus und
findet signifikante Verbesserungen vor allem bei \emph{niedrigstufigen} kognitiven Ergebnissen sowie
einen Vorteil angeleiteter gegenüber ungeleiteter Nutzung. Für höherstufige Fähigkeiten wie Erschaffen
und Bewerten fällt die Wirkung dagegen ausdrücklich minimal aus \cite{GenerativeAITools}. Genau dieses
Muster findet sich in den vorliegenden Daten wieder: Der Reflexionsbot liegt beim selbsteingeschätzten
Verständnis für Begründung und Abwägung etwas höher als der generische Chatbot, während sich auf dem
höherstufigen Zielkonstrukt der gezeigten Reflexionstiefe kein Vorteil zeigt und am objektiven
Wissenstest keine Aussage möglich ist. Der eigene Nullbefund auf der Reflexionstiefe steht damit nicht
allein, sondern entspricht dem, was die Literatur für höherstufige Kompetenzen berichtet. Die
vorliegenden Ergebnisse ergänzen diesen Befund jedoch um eine wichtige
Einschränkung: Gutes Design allein garantiert keine tiefere Reflexion, es schafft lediglich die
Gelegenheit dazu. Ein erheblicher Teil der Reflexionsbot-Gruppe nutzte diese Gelegenheit nicht,
sondern umging sie gezielt. Design und Interaktionsbereitschaft der Nutzenden sind damit als
komplementäre, nicht als austauschbare Bedingungen für den Lernerfolg zu verstehen.

Erfahrungsberichte zum Einsatz von KI-Tutoren in der Software-Engineering-Lehre ließen bislang offen,
ob solche Systeme tatsächlich zu tieferem Verständnis führen oder vor allem den Bearbeitungsaufwand
reduzieren \cite{frankfordAITutoringSoftwareEngineering2024, happeExperienceReportPedagogically2025}.
Die vorliegende Arbeit beantwortet diese Frage differenziert und mit konkreten Zahlen. Der
Reflexionsbot reduziert den Bearbeitungsaufwand nicht in dem Sinne, dass er das Ergebnis erleichtert,
sondern indem er ihn erschwert, was sich in einer geringeren Wiedernutzungsbereitschaft und einer
häufigeren Präferenz für das jeweils andere System niederschlägt
(Abschnitt~\ref{sec:beantwortung-forschungsfragen}),
und darin könnte ein Ansatzpunkt für ein tieferes Verständnis liegen. Die beiden in den
Erfahrungsberichten offengelassenen Möglichkeiten schließen sich damit nicht gegenseitig aus,
sondern hängen bei diesem Ansatz systematisch zusammen: mehr Verständnis durch mehr, nicht weniger,
Aufwand.

Auch jüngere vergleichende Studien stützen dieses Muster. Ein randomisiertes Experiment in der
Programmierausbildung findet, dass ein zurückhaltender, scaffoldender Tutor gegenüber einem direkt
antwortenden System die unmittelbare Leistung nicht garantiert steigert und kurzfristige Leistung und
tatsächliches Lernen auseinanderfallen können \cite{bassnerLessStressBetter2026}. Ebenso wird davor
gewarnt, dass müheloser KI-Einsatz eine Illusion des Lernens erzeugt, obwohl gerade die produktive
Anstrengung den Lernwert ausmacht \cite{brcicEffortlessTrapProductive2026}. Der in dieser Arbeit
beobachtete Befund, dass der Reflexionsbot trotz positiverer Einschätzung der eigenen
Kompetenzentwicklung seltener erneut genutzt werden würde und zwischen den Systemgruppen keinen
objektiven Testvorteil zeigt, fügt sich damit in ein breiteres Bild ein.

Ausdrücklich zu benennen ist jedoch eine Studie, die zu einem anderen Ergebnis kommt. Xi et al.
\cite{xiInvestigatingEffectsLLMbased2026} vergleichen einen sokratisch fragenden mit einem direkt
antwortenden Dialogagenten bei einer offenen Gestaltungsaufgabe und finden einen Vorteil der
sokratischen Bedingung gerade dort, wo die vorliegende Arbeit keinen findet, nämlich im reflexiven
Denken einschließlich der Dimension \emph{critical reflection}. Dieser Gegenbefund lässt sich nicht
mit der Aufgabenart erklären, da auch dort keine eindeutig richtige Lösung existierte. Drei
Unterschiede kommen als Erklärung eher infrage. Erstens der Umfang: Xi et al. intervenierten über
sechs Wochen mit wöchentlichen Sitzungen, die vorliegende Untersuchung über eine einzelne
Übungseinheit. Dass gerade sehr kurze Interventionen ein systematisch anderes Effektmuster zeigen,
ist in der Literatur beschrieben \cite{alemdagEffectChatbotsLearning2023}. Zweitens die Fallzahl:
94 zufällig zugeteilte Personen gegenüber 11 auswertbaren Bearbeitungen hier. Drittens die
Vergleichsbedingung: Dort stand ein eigens gebauter, in allen übrigen Eigenschaften identischer
Kontrollagent zur Verfügung, hier ein frei gewähltes Consumer-System. Die vorliegenden Daten
sprechen daher nicht gegen den Befund von Xi et al., sie zeigen vielmehr, dass sich ein solcher
Effekt unter den Bedingungen einer einzelnen Übungseinheit nicht reproduzieren lässt. Der Befund
dieser Arbeit ist damit eher als Aussage über die Grenzen kurzer Einzelinterventionen zu lesen denn
als Aussage über das Prinzip des sokratischen Fragens.

\section{Limitationen}
\label{sec:limitationen}

Die Kodierung der Reflexionsstufen und der Autorenschaft erfolgte durch eine einzelne Person. Eine
Zweitkodierung zur Berechnung einer Intercoder-Reliabilität, wie sie für strukturierende qualitative
Inhaltsanalysen üblich ist, wurde im Rahmen dieser Arbeit nicht durchgeführt. Die Einordnung
einzelner Fälle, insbesondere bei der Kategorie \glqq unklar\grqq{}, beruht damit auf einer
Einzeleinschätzung anhand von Indizien (Kürze, fehlender Chatlog, auffällige stilistische Ähnlichkeit
zwischen Personen) und nicht auf einem durch zwei unabhängige Personen abgesicherten Urteil.

\newpage
Hinzu kommt die geringe Stichprobengröße. Von 31 kodierten Fällen ist die Autorenschaft nur bei 11
eindeutig studentisch (3 bei Team AI, 8 beim Reflexionsbot). Aussagen zur Verteilung der
Reflexionsstufen, insbesondere für Team AI, sind bei dieser Fallzahl deskriptiv zu lesen und nicht
statistisch verallgemeinerbar. Für Sarma (Team AI) liegt zudem keine eigene Ausgangsantwort vor, was
im Datensatz als Ausfall dokumentiert, aber nicht ersetzt werden konnte.

Einschränkend ist zudem zu beachten, dass die zentrale Auslagerungskennzahl (80\,\% gegenüber
50\,\%) nicht unabhängig vom Aufgabendesign ist: Team AI war durch die letzte Teilaufgabe (\glqq AI
Analysis\grqq{}, vgl. \hyperref[app:szenarien]{Anhang~A.4}) ausdrücklich verpflichtet, die eigenen Entscheidungen in ein
generisches KI-System einzugeben, während die Reflexionsbot-Gruppe keine vergleichbare Vorgabe
erhielt. Ein Teil der bei Team AI beobachteten KI-Artefakte ist damit durch die Aufgabenstellung
selbst mit angelegt und nicht allein Ausdruck freiwilliger Auslagerung. Betrachtet man ausschließlich
die \emph{gesichert} KI-generierten Fälle, fällt der Gruppenunterschied deutlich geringer aus (Team
AI 8 von 15, Reflexionsbot 7 von 16, vgl. Abschnitt~\ref{sec:ergebnisse}). Der größere Abstand in der
80/50-Kennzahl entsteht erst durch die zusätzlich einbezogene, unsicher eingestufte Restkategorie
\glqq unklar\grqq{}. Die Kennzahl ist daher als Indiz, nicht als exakte Effektgröße zu lesen.

Hinzu kommt eine Asymmetrie in der Datenlage selbst, die genau diese Restkategorie betrifft. Der
Reflexionsbot stellte einen Export des vollständigen Gesprächsverlaufs bereit
(Abschnitt~\ref{sec:anforderungen}), während die Vergleichsgruppe ihre Verläufe manuell sichern
musste. Da das Vorliegen eines Chatlogs in der Autorenschaftskodierung als Indikator für studentische
Bearbeitung und sein Fehlen als Indikator für die Kategorie \glqq unklar\grqq{} eingeht
(\hyperref[app:kategoriensystem]{Anhang~A.1}), ist die Einstufung \glqq unklar\grqq{} in der
Vergleichsgruppe strukturell wahrscheinlicher. Tatsächlich entfallen vier der fünf
\glqq unklar\grqq{}-Fälle auf Team AI, und in zwei davon nennt die Kodieranmerkung das fehlende
Chatlog ausdrücklich als Grund (\hyperref[app:kodierungstabelle]{Anhang~A.2}). Der Vergleich der
gesichert KI-generierten Fälle ist aus diesem Grund die belastbarere der beiden Kennzahlen.

Eine weitere Einschränkung betrifft die Interventionstreue. Die unabhängige Variable dieser Arbeit ist
ein System, das definitionsgemäß keine direkten Lösungen liefert, sondern konsequent zurückfragt. Die
Kodierung dokumentiert jedoch drei Fälle, in denen der Bot diesem Prinzip unter anhaltendem Druck der
Studierenden nicht mehr folgte und schließlich doch direkt antwortete
(\hyperref[app:kodierungstabelle]{Anhang~A.2}, vgl. Abschnitt~\ref{sec:ausblick}). Ein systematischer
Manipulation Check, der geprüft hätte, ob die Intervention in allen Fällen wie spezifiziert wirksam
war, wurde nicht durchgeführt. Verglichen wurde damit strenggenommen nicht ein durchgängig
sokratisches System mit einem generischen, sondern ein sokratisches System, das in einem Teil der
Fälle nachgab. Da dieses Nachgeben in Richtung der Vergleichsbedingung wirkt, dürfte es die
beobachteten Gruppenunterschiede eher verkleinert als vergrößert haben; belegen lässt sich das mit den
vorliegenden Daten jedoch nicht.

Schließlich ist eine sprachliche Asymmetrie zwischen den Bedingungen zu berücksichtigen. Der
Reflexionsbot antwortet durchgängig auf Englisch, unabhängig von der Eingabesprache
(Abschnitt~\ref{sec:konzept-reflexionsbot}), während die frei gewählten Vergleichssysteme in der
Sprache der Eingabe antworten. Für überwiegend deutschsprachige Studierende im ersten Fachsemester
entsteht dadurch eine zusätzliche Verarbeitungshürde, die ausschließlich die Interventionsgruppe
betrifft. Ein Teil der beobachteten kurzen oder ausweichenden Antworten in dieser Gruppe könnte auf
die Sprachbarriere und nicht auf mangelnde Reflexionsbereitschaft zurückgehen. Die didaktische
Begründung für die englische Interaktionssprache bleibt davon unberührt, sie ersetzt aber nicht die
methodische Einschränkung.

Auch die Generalisierbarkeit der Ergebnisse ist begrenzt. Die Stichprobe stammt aus einer einzigen
Übungseinheit einer einzigen Lehrveranstaltung an einer einzigen Hochschule und besteht überwiegend
aus Studierenden des ersten Fachsemesters (vgl. Abschnitt~\ref{sec:stichprobe}). Ob sich die
beobachteten Effekte auf andere Kohorten, Studiengänge oder Hochschulen übertragen lassen, kann diese
Arbeit nicht beantworten. Die Zuteilung zu den beiden Gruppen erfolgte durch die Übungsleitung und
nicht dokumentiert randomisiert. Systematische Ausgangsunterschiede zwischen den Gruppen lassen sich
daher nicht vollständig ausschließen, auch wenn die in \hyperref[app:statistik]{Anhang~A.6} berichteten Vorerfahrungswerte auf
weitgehend vergleichbare Gruppen hindeuten. Da die Teilnahme freiwillig war, ist zudem ein
Selbstselektionseffekt nicht auszuschließen, denn Studierende, die sich für eine Teilnahme entschieden, unterscheiden sich möglicherweise
bereits in ihrer Reflexionsbereitschaft von jenen, die nicht teilnahmen. Schließlich wurde der
Reflexionsbot ausschließlich mit einem einzigen Sprachmodell (\texttt{gpt-4o-mini}) umgesetzt und
getestet. Inwiefern sich die Ergebnisse mit einem anderen oder leistungsfähigeren Modell reproduzieren
lassen, bleibt offen. Hinzu kommt eine Asymmetrie zwischen den Bedingungen. Während der Reflexionsbot
über ein festes Modell und eine feste Systemnachricht exakt spezifiziert ist, war die
Vergleichsbedingung ein frei gewähltes Consumer-System mit zum Erhebungszeitpunkt nicht
kontrolliertem Modell (vgl. Abschnitt~\ref{sec:fragebogen-experiment}). Der Vergleich ist daher als
Kontrast zwischen einem gezielt gestalteten und einem typischen alltäglichen KI-Zugang zu verstehen,
nicht als kontrollierter Modellvergleich. Schließlich wurde nicht systematisch geprüft, ob der
Reflexionsbot in Einzelfällen sachlich falsche Rückfragen stellte (Halluzinationen, vgl.
Abschnitt~\ref{sec:relevante-technologien-und-konzepte}). Inhaltlich fehlgeleitete Rückfragen könnten
die Reflexion vereinzelt in eine falsche Richtung gelenkt haben.

\newpage
\section{Handlungsempfehlungen für den Einsatz in der Software-Engineering-Lehre}
\label{sec:handlungsempfehlungen}

Aus den berichteten Befunden lassen sich sechs Empfehlungen für Lehrende ableiten, die ein
reflexionsorientiertes KI-System in einer Lehrveranstaltung einsetzen möchten. Sie sind als
Gestaltungshinweise für vergleichbare Übungssettings zu verstehen und teilen die Reichweite der in
Abschnitt~\ref{sec:limitationen} genannten Einschränkungen, insbesondere die kleine Stichprobe und den
einmaligen Erhebungszeitpunkt.

\textbf{Das Werkzeug in Aufgaben- und Bewertungsdesign einbetten, nicht danebenstellen.} Der zentrale
Befund dieser Arbeit ist, dass das Systemdesign allein die Auseinandersetzung nicht erzwingt. Ein
erheblicher Teil der Reflexionsbot-Gruppe wich den Rückfragen aus oder versuchte gezielt, das System
zu umgehen (Abschnitt~\ref{sec:beantwortung-forschungsfragen}). Ein Reflexionswerkzeug sollte deshalb
nicht als optionales Zusatzangebot eingeführt werden, sondern in eine Aufgabenstellung eingebettet
sein, deren Bearbeitung ohne die Auseinandersetzung mit den Rückfragen nicht sinnvoll möglich ist.

\textbf{Den Prozess bewerten, nicht das Produkt.} Die Kodierung zeigt, dass sprachlich und strukturell
überzeugende Abgaben in dieser Stichprobe ein stärkeres Indiz für KI-Autorenschaft waren als für tiefe
eigene Reflexion (Abschnitt~\ref{sec:illustration-kategoriensystem}). Ein fertiger Text ist damit als
Beurteilungsgrundlage für Reflexionsleistung nur noch eingeschränkt geeignet. Der Gesprächsverlauf mit
dem System ist es dagegen sehr wohl, da in ihm sichtbar wird, ob eine Person auf Rückfragen eingegangen
ist. Lehrende sollten den Chatverlauf daher zum verpflichtenden Bestandteil der Abgabe machen und ihn
in die Bewertung einbeziehen.

\textbf{Das System gegen Umgehungsversuche härten.} In mehreren dokumentierten Fällen gab der Bot
unter anhaltendem Druck nach und lieferte doch eine direkte Antwort
(\hyperref[app:kodierungstabelle]{Anhang~A.2}). Wer ein solches System einsetzt, sollte damit rechnen,
dass Studierende diese Grenze aktiv testen, und die Systemnachricht entsprechend absichern, etwa durch
eine explizite Regel für wiederholte Direktanfragen sowie durch Erkennung eingefügter vollständiger
Aufgabenstellungen (vgl. Abschnitt~\ref{sec:ausblick}).

\textbf{Erwartungen vorab rahmen.} Der Reflexionsbot übertraf die Erwartungen der Studierenden
deutlich stärker als das generische System, wurde zugleich aber seltener zur erneuten Nutzung gewählt
(Abschnitt~\ref{sec:ergebnisse}). Ein System, das bewusst keine Lösungen liefert, widerspricht der
Nutzungserfahrung, die Studierende aus generischen Assistenzsystemen mitbringen. Eine kurze
Einordnung vor dem Einsatz, warum das System nicht antwortet und worin der Lernzweck dieser
Zurückhaltung besteht, dürfte die Akzeptanz erhöhen und ist mit geringem Aufwand umsetzbar.

\textbf{Aufgabentypen gezielt auswählen.} Die vorliegende Arbeit setzt sokratisches Nachfragen bei
Prozessentscheidungen ohne eindeutig richtige Lösung ein, also bei offenen Dilemmata mit konkurrierenden
Zielkriterien. Für solche Aufgaben erwies sich die Szenario-Spezifität der Rückfragen als Voraussetzung
dafür, dass die adaptive Tiefe präzise auf einzelne Argumentationslücken reagieren konnte
(Abschnitt~\ref{sec:einordnung-forschungsstand}). Bei Aufgaben mit eindeutiger Lösung sowie bei rein
faktenbezogenen Inhalten ist der Mehrwert dagegen fraglich, zumal die Literatur die Wirkung
generativer Systeme gerade bei niedrigstufigen kognitiven Zielen als am deutlichsten beschreibt
\cite{GenerativeAITools}.

\textbf{Sprache und Zugänglichkeit bewusst entscheiden.} Die durchgängig englische Interaktionssprache
war didaktisch begründet (Abschnitt~\ref{sec:konzept-reflexionsbot}), stellt für deutschsprachige
Erstsemester aber eine zusätzliche Hürde dar, die im Vergleichssystem nicht bestand
(Abschnitt~\ref{sec:limitationen}). In einer Lehrveranstaltung ohne internationale Teilnehmende
spricht wenig dagegen, das System in der Unterrichtssprache antworten zu lassen. Andernfalls sollte
die Sprachwahl gegenüber den Studierenden begründet werden.

\section{Ausblick}
\label{sec:ausblick}

Aus den Ergebnissen und Limitationen dieser Arbeit lassen sich drei Richtungen für die
Weiterentwicklung ableiten: die technische Härtung des Reflexionsbots selbst, die methodische
Absicherung der Auswertung sowie offene Forschungsfragen, die über den Rahmen dieser Arbeit
hinausgehen.

\textbf{Technische Weiterentwicklung des Reflexionsbots.} Mehrere Fälle in der vorliegenden
Stichprobe zeigen, dass Studierende gezielt versuchten, die Rückfrage-Logik des Bots zu umgehen: Eine
Person forderte wiederholt \glqq Alternative Lösung bitte\grqq{} und drohte dem System, um eine
direkte Lösung zu erzwingen, eine andere drängte so lange auf Direktantworten, bis der Bot nachgab,
und eine dritte legte dem Bot einen fertigen Lösungsansatz vor, als kenne sie das Thema bereits, um
ihn zu einer Bestätigung statt einer Rückfrage zu bewegen. In allen drei Fällen gab der Bot dem Druck
letztlich nach. Eine Weiterentwicklung sollte deshalb gezielt gegen solche Umgehungsversuche
gehärtet werden, etwa indem copy-paste-artige Eingaben (vollständige Aufgabenstellungen oder bereits
fertige Lösungen) erkannt und mit einer Rückfrage statt einer Bestätigung beantwortet werden, und
indem die Systemnachricht explizit anweist, auf wiederholten Druck nicht nachzugeben, sondern die
Rückfrage-Strategie konsequent beizubehalten. Prompt-Engineering-Muster für robuste,
manipulationsresistente Systemnachrichten \cite{whitePromptPatternCatalog2023} bieten dafür einen
methodischen Ausgangspunkt.

\textbf{Methodische Weiterentwicklung.} Wie in Abschnitt~\ref{sec:limitationen} dargelegt, wurde die
Kodierung in dieser Arbeit von einer einzelnen Person vorgenommen. Eine Folgeuntersuchung sollte eine
echte Zweitkodierung durch eine unabhängige Person vorsehen und die Übereinstimmung, etwa über Cohens
Kappa, berichten. Ebenso sollte das Mini-Experiment über mehrere Semester oder Kohorten wiederholt
werden, um die aktuell sehr kleine studentische Teilstichprobe (n=11 von 31 Fällen) zu vergrößern und
belastbarere Aussagen zur Verteilung der Reflexionsstufen je Gruppe zu ermöglichen. Eine Ausweitung
auf weitere Lehrveranstaltungen, Studiengänge oder Hochschulen würde zugleich prüfbar machen, ob sich
die Befunde über den hier untersuchten Kontext hinaus übertragen lassen.

Am Untersuchungsdesign selbst lassen sich drei weitere Einschränkungen gezielt beheben. Beide Gruppen
sollten identische Aufgabenanforderungen erhalten, damit die Auslagerungskennzahl nicht durch die
Aufgabenstellung selbst mitbestimmt wird. Der Gegencheck mit einem generischen System müsste dafür
entweder für beide Gruppen vorgesehen sein oder für beide entfallen. Die Zuteilung zu den Gruppen
sollte zudem dokumentiert randomisiert erfolgen, um systematische Ausgangsunterschiede auszuschließen.
Um den Selbstselektionseffekt zu verringern, bietet sich statt einer freiwilligen Anmeldung eine
Vollerhebung innerhalb der regulären Übung mit Widerspruchsmöglichkeit an.

Auch die Vergleichsbedingung und die eingesetzten Instrumente lassen sich schärfen. Statt eines frei
gewählten Consumer-Systems könnte die Vergleichsgruppe dasselbe Sprachmodell mit einer neutralen
Systemnachricht nutzen. Damit ließe sich der didaktische Gestaltungseffekt vom Modelleffekt trennen,
der im vorliegenden Design vermischt bleibt. Eine Replikation mit mindestens einem weiteren
Sprachmodell würde zusätzlich zeigen, ob die Befunde modellabhängig sind. Der Wissenstest sollte
schwieriger angelegt werden, da beide Gruppen bereits im Prä-Test nahe der Höchstpunktzahl lagen und
sich Lernzuwächse dadurch kaum abbilden ließen (vgl. Abschnitt~\ref{sec:ergebnisse}). Schließlich
sollten die Chatverläufe zusätzlich daraufhin kodiert werden, ob die Rückfragen des Bots sachlich
zutreffend waren, um mögliche Fehlleitungen durch Halluzinationen von den Effekten des
Rückfrage-Prinzips zu trennen.

\textbf{Offene Forschungsfrage.} Die Ergebnisse dieser Arbeit legen einen Zielkonflikt offen, der
über den Rahmen dieser Arbeit hinausgeht: Der Reflexionsbot steht zwar bei den Studierenden, die
sich auf ihn einlassen, mit tieferer Reflexion in Verbindung, wird aber gleichzeitig als anstrengender
empfunden und seltener zur erneuten Nutzung gewählt als ein generisches System
(Abschnitt~\ref{sec:beantwortung-forschungsfragen}). Künftige Arbeiten könnten systematisch
untersuchen, wie viel Widerstand durch Rückfragen optimal ist, bevor Studierende die
Auseinandersetzung ganz vermeiden, und ob sich dieser Zielkonflikt durch ein noch stärker adaptives
Vorgehen entschärfen lässt, das die Hartnäckigkeit der Rückfragen individuell an die
Frustrationstoleranz der jeweiligen Person anpasst.
